% Kuvasarja: for n in 1 2 3 4; do echo "\\def\\step{$n}\% Kuvasarja: for n in 1 2 3 4; do echo "\\def\\step{$n}\% Kuvasarja: for n in 1 2 3 4; do echo "\\def\\step{$n}\% Kuvasarja: for n in 1 2 3 4; do echo "\\def\\step{$n}\\input{mutaatio.tex}" > mutaatio-$n.tex; tectonic mutaatio-$n.tex; rm mutaatio-$n.tex; done
\providecommand\step{4}
\documentclass[tikz,border=0pt]{standalone}
\usepackage[T1]{fontenc}
\usepackage[utf8]{inputenc}
\usepackage[scaled]{helvet}\renewcommand\familydefault{\sfdefault}
\usepackage[scaled=0.92]{beramono}
\usetikzlibrary{arrows.meta,positioning,backgrounds,calc}
\def\mkstep#1{\ifnum\step<#1\relax\tikzset{s#1/.style={opacity=0}}\else\tikzset{s#1/.style={}}\fi}
\mkstep{2}\mkstep{3}\mkstep{4}
\newcommand\hl[1]{\colorbox{orange!35}{\strut #1}}
\begin{document}
\begin{tikzpicture}[show background rectangle, background rectangle/.style={fill=white}, inner frame sep=6pt,
  code/.style={draw, thick, rounded corners=3pt, fill=#1, align=left, font=\ttfamily\small, inner sep=6pt, text width=7.2cm},
  code/.default=gray!8,
  ttl/.style={font=\small\bfseries, anchor=south west, inner sep=2pt},
  arr/.style={-{Stealth[length=3mm]}, thick},
  ok/.style={draw, thick, rounded corners=3pt, fill=green!20, font=\small\bfseries, align=center, text width=2.6cm, minimum height=1cm},
  bad/.style={draw, thick, rounded corners=3pt, fill=red!20, font=\small\bfseries, align=center, text width=2.6cm, minimum height=1cm}]

% alkuperainen koodi ja testit
\node[code] (orig) at (0,0) {def ikaluokka(ika):\\
\ \ \ \ if ika < 18:\\
\ \ \ \ \ \ \ \ return "alaikäinen"\\
\ \ \ \ return "aikuinen"};
\node[ttl] at (orig.north west) {Koodi};
\node[code=blue!6, below=0.9cm of orig] (tests) {assert ikaluokka(10) == "alaikäinen"\\
assert ikaluokka(40) == "aikuinen"};
\node[ttl] at (tests.north west) {Testit};
\node[font=\footnotesize\itshape, below=0.12cm of tests] {rivikattavuus 100\,\%};

% mutantit
\begin{scope}[s2]
\node[code=orange!6, text width=3.6cm, right=1.3cm of orig.north east, anchor=north west, yshift=0.35cm] (m1) {if ika \hl{>} 18:};
\node[ttl] at (m1.north west) {Mutantti 1};
\node[code=orange!6, text width=3.6cm, below=0.6cm of m1] (m2) {if ika \hl{<=} 18:};
\node[ttl] at (m2.north west) {Mutantti 2};
\node[code=orange!6, text width=3.6cm, below=0.6cm of m2] (m3) {return \hl{None}};
\node[ttl] at (m3.north west) {Mutantti 3};
\draw[arr] (orig.east) -- ++(0.7,0) |- (m1.west);
\draw[arr] (orig.east) ++(0.7,0) |- (m2.west);
\draw[arr] (orig.east) ++(0.7,0) |- (m3.west);
\node[font=\footnotesize\itshape, align=center, above=0.05cm] at ($(orig.east)+(0.75,0)$) {};
\end{scope}

% tulokset
\begin{scope}[s3]
\node[ok, right=0.6cm of m1] (r1) {testi hajoaa\\{\normalfont tapettu}};
\node[bad, right=0.6cm of m2] (r2) {testit menevät läpi\\{\normalfont elää!}};
\node[ok, right=0.6cm of m3] (r3) {testi hajoaa\\{\normalfont tapettu}};
\draw[arr] (m1) -- (r1); \draw[arr] (m2) -- (r2); \draw[arr] (m3) -- (r3);
\node[font=\small\bfseries, align=left, anchor=north west] at ($(m3.south west)+(0,-0.35)$) {mutaatiopisteet 2/3 = 67\,\%};
\end{scope}

% korjaus
\begin{scope}[s4]
\node[code=green!10, text width=7.2cm, anchor=north west] (fix) at ($(tests.south west)+(0,-1.25)$) {assert ikaluokka(18) == "aikuinen"};
\node[ttl] at (fix.north west) {Puuttuva raja-arvotesti, tappaa mutantin 2};
\end{scope}
\end{tikzpicture}
\end{document}
" > mutaatio-$n.tex; tectonic mutaatio-$n.tex; rm mutaatio-$n.tex; done
\providecommand\step{4}
\documentclass[tikz,border=0pt]{standalone}
\usepackage[T1]{fontenc}
\usepackage[utf8]{inputenc}
\usepackage[scaled]{helvet}\renewcommand\familydefault{\sfdefault}
\usepackage[scaled=0.92]{beramono}
\usetikzlibrary{arrows.meta,positioning,backgrounds,calc}
\def\mkstep#1{\ifnum\step<#1\relax\tikzset{s#1/.style={opacity=0}}\else\tikzset{s#1/.style={}}\fi}
\mkstep{2}\mkstep{3}\mkstep{4}
\newcommand\hl[1]{\colorbox{orange!35}{\strut #1}}
\begin{document}
\begin{tikzpicture}[show background rectangle, background rectangle/.style={fill=white}, inner frame sep=6pt,
  code/.style={draw, thick, rounded corners=3pt, fill=#1, align=left, font=\ttfamily\small, inner sep=6pt, text width=7.2cm},
  code/.default=gray!8,
  ttl/.style={font=\small\bfseries, anchor=south west, inner sep=2pt},
  arr/.style={-{Stealth[length=3mm]}, thick},
  ok/.style={draw, thick, rounded corners=3pt, fill=green!20, font=\small\bfseries, align=center, text width=2.6cm, minimum height=1cm},
  bad/.style={draw, thick, rounded corners=3pt, fill=red!20, font=\small\bfseries, align=center, text width=2.6cm, minimum height=1cm}]

% alkuperainen koodi ja testit
\node[code] (orig) at (0,0) {def ikaluokka(ika):\\
\ \ \ \ if ika < 18:\\
\ \ \ \ \ \ \ \ return "alaikäinen"\\
\ \ \ \ return "aikuinen"};
\node[ttl] at (orig.north west) {Koodi};
\node[code=blue!6, below=0.9cm of orig] (tests) {assert ikaluokka(10) == "alaikäinen"\\
assert ikaluokka(40) == "aikuinen"};
\node[ttl] at (tests.north west) {Testit};
\node[font=\footnotesize\itshape, below=0.12cm of tests] {rivikattavuus 100\,\%};

% mutantit
\begin{scope}[s2]
\node[code=orange!6, text width=3.6cm, right=1.3cm of orig.north east, anchor=north west, yshift=0.35cm] (m1) {if ika \hl{>} 18:};
\node[ttl] at (m1.north west) {Mutantti 1};
\node[code=orange!6, text width=3.6cm, below=0.6cm of m1] (m2) {if ika \hl{<=} 18:};
\node[ttl] at (m2.north west) {Mutantti 2};
\node[code=orange!6, text width=3.6cm, below=0.6cm of m2] (m3) {return \hl{None}};
\node[ttl] at (m3.north west) {Mutantti 3};
\draw[arr] (orig.east) -- ++(0.7,0) |- (m1.west);
\draw[arr] (orig.east) ++(0.7,0) |- (m2.west);
\draw[arr] (orig.east) ++(0.7,0) |- (m3.west);
\node[font=\footnotesize\itshape, align=center, above=0.05cm] at ($(orig.east)+(0.75,0)$) {};
\end{scope}

% tulokset
\begin{scope}[s3]
\node[ok, right=0.6cm of m1] (r1) {testi hajoaa\\{\normalfont tapettu}};
\node[bad, right=0.6cm of m2] (r2) {testit menevät läpi\\{\normalfont elää!}};
\node[ok, right=0.6cm of m3] (r3) {testi hajoaa\\{\normalfont tapettu}};
\draw[arr] (m1) -- (r1); \draw[arr] (m2) -- (r2); \draw[arr] (m3) -- (r3);
\node[font=\small\bfseries, align=left, anchor=north west] at ($(m3.south west)+(0,-0.35)$) {mutaatiopisteet 2/3 = 67\,\%};
\end{scope}

% korjaus
\begin{scope}[s4]
\node[code=green!10, text width=7.2cm, anchor=north west] (fix) at ($(tests.south west)+(0,-1.25)$) {assert ikaluokka(18) == "aikuinen"};
\node[ttl] at (fix.north west) {Puuttuva raja-arvotesti, tappaa mutantin 2};
\end{scope}
\end{tikzpicture}
\end{document}
" > mutaatio-$n.tex; tectonic mutaatio-$n.tex; rm mutaatio-$n.tex; done
\providecommand\step{4}
\documentclass[tikz,border=0pt]{standalone}
\usepackage[T1]{fontenc}
\usepackage[utf8]{inputenc}
\usepackage[scaled]{helvet}\renewcommand\familydefault{\sfdefault}
\usepackage[scaled=0.92]{beramono}
\usetikzlibrary{arrows.meta,positioning,backgrounds,calc}
\def\mkstep#1{\ifnum\step<#1\relax\tikzset{s#1/.style={opacity=0}}\else\tikzset{s#1/.style={}}\fi}
\mkstep{2}\mkstep{3}\mkstep{4}
\newcommand\hl[1]{\colorbox{orange!35}{\strut #1}}
\begin{document}
\begin{tikzpicture}[show background rectangle, background rectangle/.style={fill=white}, inner frame sep=6pt,
  code/.style={draw, thick, rounded corners=3pt, fill=#1, align=left, font=\ttfamily\small, inner sep=6pt, text width=7.2cm},
  code/.default=gray!8,
  ttl/.style={font=\small\bfseries, anchor=south west, inner sep=2pt},
  arr/.style={-{Stealth[length=3mm]}, thick},
  ok/.style={draw, thick, rounded corners=3pt, fill=green!20, font=\small\bfseries, align=center, text width=2.6cm, minimum height=1cm},
  bad/.style={draw, thick, rounded corners=3pt, fill=red!20, font=\small\bfseries, align=center, text width=2.6cm, minimum height=1cm}]

% alkuperainen koodi ja testit
\node[code] (orig) at (0,0) {def ikaluokka(ika):\\
\ \ \ \ if ika < 18:\\
\ \ \ \ \ \ \ \ return "alaikäinen"\\
\ \ \ \ return "aikuinen"};
\node[ttl] at (orig.north west) {Koodi};
\node[code=blue!6, below=0.9cm of orig] (tests) {assert ikaluokka(10) == "alaikäinen"\\
assert ikaluokka(40) == "aikuinen"};
\node[ttl] at (tests.north west) {Testit};
\node[font=\footnotesize\itshape, below=0.12cm of tests] {rivikattavuus 100\,\%};

% mutantit
\begin{scope}[s2]
\node[code=orange!6, text width=3.6cm, right=1.3cm of orig.north east, anchor=north west, yshift=0.35cm] (m1) {if ika \hl{>} 18:};
\node[ttl] at (m1.north west) {Mutantti 1};
\node[code=orange!6, text width=3.6cm, below=0.6cm of m1] (m2) {if ika \hl{<=} 18:};
\node[ttl] at (m2.north west) {Mutantti 2};
\node[code=orange!6, text width=3.6cm, below=0.6cm of m2] (m3) {return \hl{None}};
\node[ttl] at (m3.north west) {Mutantti 3};
\draw[arr] (orig.east) -- ++(0.7,0) |- (m1.west);
\draw[arr] (orig.east) ++(0.7,0) |- (m2.west);
\draw[arr] (orig.east) ++(0.7,0) |- (m3.west);
\node[font=\footnotesize\itshape, align=center, above=0.05cm] at ($(orig.east)+(0.75,0)$) {};
\end{scope}

% tulokset
\begin{scope}[s3]
\node[ok, right=0.6cm of m1] (r1) {testi hajoaa\\{\normalfont tapettu}};
\node[bad, right=0.6cm of m2] (r2) {testit menevät läpi\\{\normalfont elää!}};
\node[ok, right=0.6cm of m3] (r3) {testi hajoaa\\{\normalfont tapettu}};
\draw[arr] (m1) -- (r1); \draw[arr] (m2) -- (r2); \draw[arr] (m3) -- (r3);
\node[font=\small\bfseries, align=left, anchor=north west] at ($(m3.south west)+(0,-0.35)$) {mutaatiopisteet 2/3 = 67\,\%};
\end{scope}

% korjaus
\begin{scope}[s4]
\node[code=green!10, text width=7.2cm, anchor=north west] (fix) at ($(tests.south west)+(0,-1.25)$) {assert ikaluokka(18) == "aikuinen"};
\node[ttl] at (fix.north west) {Puuttuva raja-arvotesti, tappaa mutantin 2};
\end{scope}
\end{tikzpicture}
\end{document}
" > mutaatio-$n.tex; tectonic mutaatio-$n.tex; rm mutaatio-$n.tex; done
\providecommand\step{4}
\documentclass[tikz,border=0pt]{standalone}
\usepackage[T1]{fontenc}
\usepackage[utf8]{inputenc}
\usepackage[scaled]{helvet}\renewcommand\familydefault{\sfdefault}
\usepackage[scaled=0.92]{beramono}
\usetikzlibrary{arrows.meta,positioning,backgrounds,calc}
\def\mkstep#1{\ifnum\step<#1\relax\tikzset{s#1/.style={opacity=0}}\else\tikzset{s#1/.style={}}\fi}
\mkstep{2}\mkstep{3}\mkstep{4}
\newcommand\hl[1]{\colorbox{orange!35}{\strut #1}}
\begin{document}
\begin{tikzpicture}[show background rectangle, background rectangle/.style={fill=white}, inner frame sep=6pt,
  code/.style={draw, thick, rounded corners=3pt, fill=#1, align=left, font=\ttfamily\small, inner sep=6pt, text width=7.2cm},
  code/.default=gray!8,
  ttl/.style={font=\small\bfseries, anchor=south west, inner sep=2pt},
  arr/.style={-{Stealth[length=3mm]}, thick},
  ok/.style={draw, thick, rounded corners=3pt, fill=green!20, font=\small\bfseries, align=center, text width=2.6cm, minimum height=1cm},
  bad/.style={draw, thick, rounded corners=3pt, fill=red!20, font=\small\bfseries, align=center, text width=2.6cm, minimum height=1cm}]

% alkuperainen koodi ja testit
\node[code] (orig) at (0,0) {def ikaluokka(ika):\\
\ \ \ \ if ika < 18:\\
\ \ \ \ \ \ \ \ return "alaikäinen"\\
\ \ \ \ return "aikuinen"};
\node[ttl] at (orig.north west) {Koodi};
\node[code=blue!6, below=0.9cm of orig] (tests) {assert ikaluokka(10) == "alaikäinen"\\
assert ikaluokka(40) == "aikuinen"};
\node[ttl] at (tests.north west) {Testit};
\node[font=\footnotesize\itshape, below=0.12cm of tests] {rivikattavuus 100\,\%};

% mutantit
\begin{scope}[s2]
\node[code=orange!6, text width=3.6cm, right=1.3cm of orig.north east, anchor=north west, yshift=0.35cm] (m1) {if ika \hl{>} 18:};
\node[ttl] at (m1.north west) {Mutantti 1};
\node[code=orange!6, text width=3.6cm, below=0.6cm of m1] (m2) {if ika \hl{<=} 18:};
\node[ttl] at (m2.north west) {Mutantti 2};
\node[code=orange!6, text width=3.6cm, below=0.6cm of m2] (m3) {return \hl{None}};
\node[ttl] at (m3.north west) {Mutantti 3};
\draw[arr] (orig.east) -- ++(0.7,0) |- (m1.west);
\draw[arr] (orig.east) ++(0.7,0) |- (m2.west);
\draw[arr] (orig.east) ++(0.7,0) |- (m3.west);
\node[font=\footnotesize\itshape, align=center, above=0.05cm] at ($(orig.east)+(0.75,0)$) {};
\end{scope}

% tulokset
\begin{scope}[s3]
\node[ok, right=0.6cm of m1] (r1) {testi hajoaa\\{\normalfont tapettu}};
\node[bad, right=0.6cm of m2] (r2) {testit menevät läpi\\{\normalfont elää!}};
\node[ok, right=0.6cm of m3] (r3) {testi hajoaa\\{\normalfont tapettu}};
\draw[arr] (m1) -- (r1); \draw[arr] (m2) -- (r2); \draw[arr] (m3) -- (r3);
\node[font=\small\bfseries, align=left, anchor=north west] at ($(m3.south west)+(0,-0.35)$) {mutaatiopisteet 2/3 = 67\,\%};
\end{scope}

% korjaus
\begin{scope}[s4]
\node[code=green!10, text width=7.2cm, anchor=north west] (fix) at ($(tests.south west)+(0,-1.25)$) {assert ikaluokka(18) == "aikuinen"};
\node[ttl] at (fix.north west) {Puuttuva raja-arvotesti, tappaa mutantin 2};
\end{scope}
\end{tikzpicture}
\end{document}
